\section{Arquitectura y superficie de ataque}
\label{sec:arquitectura}

SOS Clinic es un prototipo académico organizado como un conjunto de servicios
independientes detrás de un gateway. El frontend es estático y consume las APIs
por HTTP; el gateway enruta hacia los microservicios, que se registran en Eureka.
La analítica usa una base agregada y el agente conversacional consulta las APIs de
pacientes y triage y, para la transcripción de voz, la API externa de Groq.

\begin{figure}[htbp]
\centering
\resizebox{\textwidth}{!}{%
\begin{tikzpicture}[
  node distance=8mm and 10mm,
  box/.style={draw=teal!70,fill=teal!6,rounded corners=2pt,align=center,
    minimum width=2.9cm,minimum height=0.85cm,font=\footnotesize},
  svc/.style={box,minimum width=3.0cm},
  ext/.style={box,draw=high!75,fill=high!8},
  infra/.style={draw=slategray!50,fill=slatelight,rounded corners=2pt,align=center,
    font=\footnotesize,minimum height=0.7cm,minimum width=2.2cm},
  arr/.style={-{Latex[length=1.8mm]},draw=slategray!80,thick},
  dasharr/.style={-{Latex[length=1.8mm]},draw=slategray!55,dashed}
]
  \node[box,fill=teal!12] (browser) at (0,0) {Navegador};
  \node[box] (front) at (3.4,0) {Frontend estático\\\texttt{:5500}};
  \node[box,fill=teal!15] (gw) at (7.0,0) {\texttt{ms-gateway}\\\texttt{:8080}};

  \node[svc] (seg) at (11.0,2.6) {\texttt{ms-seguridad}\\\texttt{:8081}};
  \node[svc] (pac) at (11.0,1.3) {\texttt{ms-pacientes}\\\texttt{:8082}};
  \node[svc] (tri) at (11.0,0.0) {\texttt{ms-triage}\\\texttt{:8083}};
  \node[svc] (ana) at (11.0,-1.3) {\texttt{ms-analitica}\\\texttt{:8084}};
  \node[svc] (ia)  at (11.0,-2.6) {\texttt{ms-ia-agente}\\\texttt{:8088}};

  \coordinate (bus) at (9.2,0);
  \draw[arr] (browser) -- (front);
  \draw[arr] (front) -- (gw);
  \draw[arr] (gw.east) -- (bus);
  \foreach \y in {2.6,1.3,0.0,-1.3,-2.6} {
    \draw[arr] (bus) -- (9.2,\y) -- (9.4,\y);
  }

  \node[ext] (groq) at (15.0,-2.6) {Groq API\\Transcripción};

  \node[infra] (pg) at (3.2,-4.4) {PostgreSQL\\\texttt{:5433}};
  \node[infra] (pga) at (6.3,-4.4) {pgAdmin\\\texttt{:5050}};
  \node[infra] (eureka) at (9.4,-4.4) {Eureka\\\texttt{:8761}};
  \node[infra] (mail) at (12.5,-4.4) {Mailpit\\\texttt{:8025}};

  \draw[dasharr] (ia.east) -- (groq.west);
  \draw[dasharr] (tri.south) -- ++(0,-1.05) -| (pg.north);
  \draw[dasharr] (ana.south) -- ++(0,-0.65) -| (pg.north);
  \draw[dasharr] (seg.south) -- ++(0,-0.65) -| (mail.north);
  \draw[dasharr] (gw.south) -- ++(0,-1.7) -| (eureka.north);

  \node[anchor=north west,font=\footnotesize,align=left,text=slategray]
    at (-1.0,-5.5) {Línea continua: tráfico HTTP. Línea discontinua: dependencia local o servicio externo.};
\end{tikzpicture}%
}
\caption{Arquitectura lógica y superficie de ataque de SOS Clinic.}
\label{fig:arquitectura}
\end{figure}

\subsection{Fronteras de confianza}

Las fronteras principales son navegador--frontend (controles de sesión en
JavaScript y JWT en el almacenamiento de sesión del navegador), frontend--gateway
(sin validación global del token), gateway--microservicios (pacientes y triage sin filtro),
microservicios--PostgreSQL (credenciales de desarrollo y puerto 5433) y
agente--Groq (audio a un tercero y cuota sin identidad del usuario final).

\subsection{Superficie de ataque}

\begin{itemize}
  \item \textbf{APIs de negocio sin autenticación.} \texttt{ms-pacientes} y
  \texttt{ms-triage} exponen lectura y escritura de datos clínicos sin token
  (H-02, H-03).
  \item \textbf{Rutas administrativas sin autorización por rol.} Cualquier token
  válido, incluido el de un usuario de triage, accede a \texttt{/admin/**}
  (H-04).
  \item \textbf{Secreto JWT compartido y versionado.} El valor por defecto permite
  forjar tokens de administrador (H-01).
  \item \textbf{Frontend que confía en el cliente.} La sesión y la navegación se
  controlan en JavaScript; el XSS almacenado puede extraer el token (H-07, H-23).
  \item \textbf{Agente IA con acceso directo a datos clínicos.} Responde sin token
  y no limita herramientas, tamaño de audio ni costo (H-10, H-21).
  \item \textbf{Infraestructura de desarrollo expuesta.} PostgreSQL, pgAdmin,
  Eureka y Mailpit escuchan en la red local con credenciales triviales (H-19).
\end{itemize}

\subsection{Modelado STRIDE/DREAD}

La tabla~\ref{tab:stride} resume la priorización cualitativa por componente. Los
niveles son una síntesis del análisis del informe fuente; no sustituyen un cálculo
formal de DREAD ni un vector CVSS.

\begin{table}[htbp]
\centering
\caption{Resumen STRIDE/DREAD por componente.}
\label{tab:stride}
\small
\begin{tabular}{@{}L{3.1cm}C{1.05cm}C{1.05cm}C{1.05cm}C{1.05cm}C{1.05cm}C{1.05cm}L{1.7cm}@{}}
\toprule
\textbf{Componente} & \textbf{S} & \textbf{T} & \textbf{R} & \textbf{I} & \textbf{D} & \textbf{E} & \textbf{Riesgo}\\
\midrule
\texttt{ms-seguridad} & Alta & Alta & Media & Alta & Media & Alta & \crit{Crítico}\\
\texttt{ms-pacientes} & Alta & Alta & Baja & Alta & Media & Media & \crit{Crítico}\\
\texttt{ms-triage} & Alta & Alta & Baja & Alta & Media & Media & \crit{Crítico}\\
\texttt{ms-gateway} & Media & Media & Baja & Media & Media & Alta & \alta{Alto}\\
\texttt{ms-analitica} & Baja & Baja & Baja & Media & Baja & Baja & \media{Medio}\\
\texttt{ms-ia-agente} & Alta & Media & Media & Alta & Media & Alta & \alta{Alto}\\
Frontend & Media & Media & Baja & Alta & Baja & Media & \alta{Alto}\\
\bottomrule
\end{tabular}
\end{table}

\noindent\footnotesize S: suplantación; T: manipulación; R: repudio; I: divulgación
de información; D: denegación de servicio; E: elevación de privilegios.
\normalsize

